% Generated draft: requires mathematical and source review.
\documentclass[11pt]{article}
\usepackage[margin=25mm]{geometry}
\usepackage{fontspec}
\setmainfont{Latin Modern Roman}
\usepackage{amsmath,amssymb,mathtools,unicode-math}
\setmathfont{Latin Modern Math}
\usepackage{enumitem,needspace,xurl,hyperref}
\setlength{\parindent}{0pt}
\setlength{\parskip}{5pt}
\setlength{\emergencystretch}{3em}
\newcommand{\N}{\mathbb{N}}
\newcommand{\Z}{\mathbb{Z}}
\newcommand{\Q}{\mathbb{Q}}
\newcommand{\R}{\mathbb{R}}
\newcommand{\C}{\mathbb{C}}
\newcommand{\F}{\mathbb{F}}
\newcommand{\A}{\mathbb{A}}
\newcommand{\PP}{\mathbb{P}}
\newcommand{\E}{\mathbb{E}}
\newcommand{\eps}{\varepsilon}
\newcommand{\dd}{\,\mathrm{d}}
\begin{document}
\section*{UnsolvedMath ID: 30004324}
\textbf{Source problem code:} catalogue entry
\subsection{Definitions and mathematical statement}
The source asks for a characterization of projective space by positivity of a tangent-bundle Seshadri constant. It does not explicitly specify the ground field, the range of dimensions, or the convention for Seshadri constants of vector bundles that need not be nef. The following formulation uses the usual complex-projective interpretation and explicitly states a curve-based convention; these choices require confirmation against the cited report.

Let \(n\geq 1\) be an integer and let \(X\) be a connected smooth projective variety over \(\C\) of dimension \(n\). Write \(T_X\) for its tangent bundle. For a vector bundle \(E\) of positive rank on \(X\), let \(p:\PP_X(E)\to X\) denote the projective bundle parametrizing one-dimensional quotients, with tautological quotient line bundle \(\mathcal O_{\PP_X(E)}(1)\). For an irreducible curve \(\Gamma\subset\PP_X(E)\) whose image \(C=p(\Gamma)\) is a curve, let \(\nu_\Gamma:\widetilde\Gamma\to\Gamma\) be its normalization and let \(d_\Gamma\) be the degree of the induced finite morphism \(\widetilde\Gamma\to\widetilde C\). Here \(\widetilde C\) is the normalization of \(C\), and \(\operatorname{mult}_x C\) is the multiplicity of \(C\) at \(x\). Use the convention
\[
\varepsilon(E;x)=\inf_{\substack{\Gamma\subset\PP_X(E)\text{ irreducible curve}\\ C=p(\Gamma),\ \dim C=1,\ x\in C}}
\frac{\deg\!\left(\nu_\Gamma^*\mathcal O_{\PP_X(E)}(1)\right)}{d_\Gamma\operatorname{mult}_x C}.
\]
No global nefness assumption on \(E\) or \(T_X\) is imposed.

Under these explicitly stated conventions, the question is whether, for every such \(n\) and \(X\),
\[
\left(\exists x_0\in X:\ \varepsilon(T_X;x_0)>0\right)
\quad\Longrightarrow\quad
X\cong\PP^n_{\C}.
\]
The point \(x_0\) is an arbitrary point whose existence is asserted; it is not required to be general or very general. The supplied record labels this question open. That label and its accompanying literature assessment are supplied curation, not independently verified mathematical conclusions.
\subsection{Short English statement}
Does positivity of the tangent bundle's Seshadri constant at even one point force a projective manifold to be projective space? The displayed formulation makes a conventional complex-projective interpretation and a vector-bundle Seshadri convention explicit; the source does not itself specify these details.
\subsection{Sources}
\begin{itemize}
\item Supplied local UnsolvedMath registry record, permanent ID \(30004324\), source problem code \texttt{catalogue entry}, title: Characterizing Projective Space by Tangent Bundle Seshadri Constants. The record belongs to the Oberwolfach Reports Open Problems collection. No individual proposer is supplied.
\item Supplied source citation: Mini-Workshop: Seshadri Constants (2020), Oberwolfach Reports, pp.\ 3271--3313, DOI \texttt{10.4171/owr/2019/53}. \url{https://doi.org/10.4171/owr/2019/53}. The background describes Oberwolfach Reports as documenting mathematical workshops and their open-problem discussions. The registry gives proposed year 2019; the citation labels the report 2020. These supplied dates have not been reconciled.
\item Supplied literature reference: Seshadri constants for vector bundles. \url{https://doi.org/10.1016/j.jpaa.2020.106559}.
\item Supplied literature reference: Seshadri constants of tangent sheaves on toric varieties. \url{https://arxiv.org/abs/2211.17172}.
\item Supplied literature reference: Oberwolfach Reports, tangent-bundle conjecture. \url{https://ems.press/content/serial-article-files/46833}.
\item Literature assessment dated exactly \texttt{2026-08-22}: the supplied curation labels the problem open and reports that the characterization is proved for Fano varieties, surfaces, general points, and toric varieties, while the unrestricted smooth-projective statement remains conjectural. This is supplied dated curation derived from the supplied OWR material, not an independently verified assessment. The precise hypotheses and time-sensitive status of these claims require independent checking.
\item The supplied statement-status label is \texttt{corrected\_verified}. Its accompanying qualification says that only formula delimiters, notation, grammar, and typography were repaired, preserving mathematical content and scope. This annotation does not establish mathematical correctness or independently verify the cited literature.
\end{itemize}
\end{document}
